\begin{singlespace}
\noindent\hspace{1.25cm}Os três fluxogramas do GNU Radio reproduziram a teoria: a BPSK mostrou dois pontos sobre o eixo I com Q nulo, a QPSK quatro pontos a distância unitária da origem e a 16-QAM dezesseis pontos em quatro níveis por eixo. Para trocar de modulação, bastou alterar o \texttt{Maximum} e a \textit{Symbol Table}. A normalização por \sq{2} e \sq{10} igualou a energia média a 1. A \num{32000}~símbolos/s, a taxa de bits subiu de \SI{32}{kbit/s} para \SI{64}{kbit/s} e \SI{128}{kbit/s} sem mais largura de banda, mas a distância mínima caiu de $1{,}414$ para $0{,}632$, exigindo cerca de \SI{7}{dB} a mais de relação sinal-ruído na 16-QAM. Assim, a escolha da modulação equilibra vazão e robustez.
\end{singlespace}
